% !TeX program = xelatex
% UTF-8. Standalone: xelatex splenomegaly.tex from this directory.
% To include in a XePersian document, define \SplenomegalyInNotebook first.
% Source: Google Drive notebook/splenomegaly, four images in filename order.
% pic01.jpg: 20261010_052811026.jpg (1_TssB1kRFtyIXs_WERHLI8xcFH2TNPo6)
% pic02.jpg: 20261010_052816638.jpg (1YNubILmX6CR_pkCGCQO8TdDYSHkDi4LC)
% pic03.jpg: 20261010_052821858.jpg (11oVVYxITgOldbgIx_o4pnjgklbdSCXEz)
% pic04.jpg: 20261010_052826980.jpg (1xzwJG74_CgzdDdijWdjH_6eaeZA9XMoI)
\ifdefined\SplenomegalyInNotebook
  \let\SplenomegalyFinish\relax
\else
  \def\SplenomegalyFinish{\end{document}}
  \documentclass[12pt,a4paper]{article}
  \usepackage[margin=22mm]{geometry}
  \usepackage{amsmath,amssymb}
  \usepackage{graphicx}
  \usepackage[normalem]{ulem}
  \usepackage{xepersian}
  \IfFontExistsTF{Amiri}{\settextfont{Amiri}}{\settextfont{DejaVu Sans}}
  \setlatintextfont{DejaVu Serif}
  \setlength{\parindent}{0pt}
  \setlength{\parskip}{5pt}
  \setlength{\emergencystretch}{3em}
  \begin{document}
\fi
\providecommand{\SplenoUnclear}[1]{\textbf{[خوانش نامطمئن: #1]}}
\providecommand{\SplenoPage}[2]{\clearpage\section*{تصویر #1}\lr{\texttt{#2}}\par}

\begin{center}
{\Large\bfseries اسپلنومگالی}\par
\lr{Splenomegaly}\par
۵۳۵
\end{center}
\textbf{یادداشت رونویسی:} متن چهار تصویر به ترتیب برگه‌ها رونویسی شده است. فلش‌ها به صورت فهرست یا رابطه بازنویسی شده‌اند. عبارت‌ها، اختصارات و اعداد منبع حفظ شده‌اند؛ خوانش‌های نامطمئن مشخص شده‌اند. اصل تصاویر، شامل حاشیه‌ها و خط‌خوردگی‌ها، در پیوست آمده است.

\SplenoPage{۱}{pic01.jpg}
\textbf{سربرگ و حاشیه:} ۵۳۵ ـ اسپلنومگالی؛ \lr{Subject:} (خالی)؛ \lr{Date: / /}؛ شمارهٔ برگه: ۱؛ نشان پایین برگه: \lr{YASHA}.
\begin{itemize}
\item در ۱۵\% نوزادان یک طحال نرم بزرگ قابل لمس است؛ ۱۰\% کودکان و ۵\% نوجوانان.
\item در اسپلنومگالی قابل لمس، طحال حداقل ۲–۳ برابر بزرگ شده است.
\item بهترین شرایط معاینه: بیمار \lr{Supine}، یا به سمت راست بیمار.
\item لبهٔ طحال اگر بیش از \lr{2 cm} زیر لبهٔ دنده لمس شود: پاتولوژیک.
\item اگر اسپلنومگالی در زمینهٔ افزایش فشار ورید پورت باشد، عروق دیلاتهٔ غیرطبیعی (کاپوت مدوزا) دیده می‌شود.
\item ارزیابی رادیولوژیک ساختار و عملکرد: \lr{Tc-99m sulfur colloid scan}.
\end{itemize}
\SplenoUnclear{واژهٔ پس از «طحال نرم» در سطر اول به صورت «بزرگ» خوانده می‌شود.}
\textbf{حاشیهٔ کنار اسکن:} «تکنسیم». \SplenoUnclear{واژهٔ کوتاه فارسی کنار \lr{Tc-99m} به صورت «تکنسیم» خوانده می‌شود.}

شرح حال کامل باید گرفته شود، خصوصاً دربارهٔ شکایات سیستمیک (تب، تعریق شبانه، ملنا، کاهش وزن).

معاینهٔ فیزیکی کامل باید انجام شود (لنفادنوپاتی، زردی، هپاتومگالی، راش، التهاب مفاصل، پتشی، اکیموز).

حتماً باید \lr{CBC} و \lr{PBS} انجام شود.

\SplenoPage{۲}{pic02.jpg}
\textbf{سربرگ و حاشیه:} ۵۳۵؛ \lr{Subject:} (خالی)؛ \lr{Date: / /}؛ شمارهٔ برگه: ۲؛ \lr{YASHA}.
\subsection*{سودواسپلنومگالی}
\begin{itemize}
\item اگر مزانتر متصل به طحال به طور غیرمعمول بلند باشد، طحال پایین می‌آید.
\end{itemize}
\SplenoUnclear{در جملهٔ بالا واژه‌های «مزانتر» و «بلند» خوانش قطعی ندارند؛ باقی جمله «متصل به طحال به طور غیرمعمول ... باشد، طحال پایین می‌آید» است.}
\begin{itemize}
\item لوب چپ کبد اگر بزرگ شود، در \lr{ddx}.
\item هر تودهٔ \lr{LUQ}، در \lr{ddx}.
\item هماتوم طحال، در \lr{ddx}.
\item کیست طحال، در \lr{ddx}:
  \begin{itemize}
  \item مادرزادی: اپیدرموئید.
  \item اکتسابی: سودوکیست (به دنبال تروما).
  \end{itemize}
\item \lr{Splenosis}، در \lr{ddx}: به دنبال پارگی طحال؛ کاشته شدن در سایر نقاط.
\item طحال فرعی، در \lr{ddx}: در ۱۵\% افراد وجود دارد.
\item پلی‌اسپلنیسم مادرزادی (سندرم).
\item نقص قلبی، \lr{left-sided organ anomaly}، آترزی بلیاری، سودواسپلنیسم.
\end{itemize}
\SplenoUnclear{واژهٔ انتهای سطر نقص قلبی به صورت «سودواسپلنیسم» خوانده می‌شود و بخشی از آن نزدیک حاشیهٔ برگه است.}

\subsection*{هیپراسپلنیسم (افزایش عملکرد طحال)}
\begin{itemize}
\item می‌تواند سایتوپنی ایجاد کند: با مکانیسم سکوستریشن.
\item معمولاً ثانویه به یک بیماری دیگر است.
\item اگر بیماری اولیه درمان شود، معمولاً برمی‌گردد.
\item در موارد شدید اسپلنکتومی انجام می‌دهیم.
\end{itemize}

\SplenoPage{۳}{pic03.jpg}
\textbf{سربرگ و حاشیه:} ۵۳۵؛ \lr{Subject:} (خالی)؛ \lr{Date: / /}؛ شمارهٔ برگه: ۳؛ \lr{YASHA}.
\subsection*{اسپلنومگالی کانژستیو}
\begin{itemize}
\item هایپراسپلنیسم در زمینهٔ انسداد ورید هپاتیک، پورتال، طحالی.
\item بیماری ویلسون.
\item گالاکتوزمی.
\item آترزی بلیاری.
\item \lr{\(\alpha_1\)} آنتی‌تریپسین دفیشنسی.
\end{itemize}
\textbf{توضیح متصل به فهرست با خط مورب:} منجر به التهاب در کبد، فیبروز و به نهایت انسداد وریدی.
\begin{itemize}
\item اگر طحال محل انسداد وریدی باشد: اسپلنکتومی می‌کنیم.
\item شانتینگ پورتوکاوال: انسداد در پورت، هپاتیک باشد.
\end{itemize}

\subsection*{کدام بدخیمی‌ها اسپلنومگالی می‌دهند؟}
\begin{itemize}
\item لوکمی (حاد، مزمن)، لنفوم، آنژیوسارکوم، ماستوسیتوز.
\item متاستاز سایر بدخیمی‌ها به طحال.
\end{itemize}

\subsection*{بیماری‌های احتقانی که اسپلنومگالی می‌دهند}
\begin{itemize}
\item \lr{HF}، سیروز یا فیبروز کبدی، سندرم بودکیاری، ترومبوز در وریدهای کبدی و پورت.
\end{itemize}

\subsection*{عفونت‌های ویروسی عامل اسپلنومگالی}
\begin{itemize}
\item \lr{HSV}، \lr{CMV}، روبلا، \lr{HAV}، \lr{HBV}، \lr{HCV}، \lr{EBV}، \lr{HIV}، \lr{HHV-6}.
\end{itemize}

\SplenoPage{۴}{pic04.jpg}
\textbf{سربرگ و حاشیه:} ۵۳۵؛ \lr{Subject:} (خالی)؛ \lr{Date: / /}؛ شمارهٔ برگه: ۴؛ \lr{YASHA}.
\subsection*{هایپرپلازی در زمینهٔ اختلالات هماتولوژیک}
\begin{itemize}
\item همولیز حاد و مزمن.
\item هموگلوبینوپاتی‌ها: \lr{SSD}، تالاسمی مینور.
\item اختلالات غشای اریتروسیت: \lr{HS}، الیپتوسیتوز، پیروپویکیلوسیتوز.
\item نقص آنزیمی اریتروسیت: \lr{G6PDD} شدید، کمبود پیرووات کیناز.
\item همولیز ایمنی: اتوایمیون، ایزوایمیون.
\item \lr{PNH}.
\item فقر آهن مزمن.
\item هماتوپوئزیس اکسترامدولاری.
\item بیماری‌های میلوپرولیفراتیو: \lr{CNL}، میلوفیبروز با متاپلازی میلوئید، \lr{PV}.
\item استئوپتروز بدخیم.
\item دریافت \lr{GCSF}.
\end{itemize}
\SplenoUnclear{اختصار کنار میلوفیبروز در دست‌نویس «\lr{CNL}» خوانده می‌شود؛ ممکن است «\lr{CML}» باشد. اختصارات «\lr{SSD}»، «\lr{G6PDD}» و «\lr{GCSF}» با املای تصویر حفظ شده‌اند.}
\SplenoUnclear{واژهٔ پس از «استئوپتروز» به صورت «بدخیم» خوانده می‌شود.}

\subsection*{آبسهٔ طحال عفونی}
استاف اورئوس، استرپ، سالمونلا.

\textbf{خط‌خوردگی پیش از استاف اورئوس:} \SplenoUnclear{واژهٔ خط‌خورده ظاهراً «بروسلا» است؛ به دلیل خط‌خوردگی خوانش قطعی نیست.}

\subsection*{غیرنئوپلاستیک و التهاب}
\begin{itemize}
\item \lr{SLE}، \lr{JIA}، \lr{SSS}، \lr{GVHD}، \lr{MCTD}، شوگران، سارکوئیدوز.
\end{itemize}
\SplenoUnclear{اختصار سوم در این فهرست به صورت «\lr{SSS}» نوشته شده است؛ بدون بسط یا اصلاح حفظ شده است.}

\clearpage
\section*{پیوست: چهار تصویر کامل منبع}
تمام سربرگ‌ها، حاشیه‌ها، نشانه‌ها و خط‌خوردگی‌های دست‌نویس در تصاویر زیر محفوظ‌اند.
% Support compilation from this directory or the repository root.
\IfFileExists{../../ped-infectious diseases/splenomegaly/images/pic01.jpg}{%
  \def\SplenoImageBase{../../ped-infectious diseases/splenomegaly/images/}}{%
  \def\SplenoImageBase{pediatrics/ped-infectious diseases/splenomegaly/images/}}
\providecommand{\SplenoOriginal}[2]{%
  \clearpage\subsection*{اصل تصویر #1: \lr{\texttt{pic#2.jpg}}}%
  \begin{center}
  \includegraphics[width=\linewidth,height=0.87\textheight,keepaspectratio]{\SplenoImageBase pic#2.jpg}
  \end{center}}
\SplenoOriginal{۱}{01}
\SplenoOriginal{۲}{02}
\SplenoOriginal{۳}{03}
\SplenoOriginal{۴}{04}
\SplenomegalyFinish
